\begin{ExerciseBlock}
\ExerciseHeading{06}{Schwingendes Reagenzglas}{free_oscillations_floating_test_tube}
Ein Reagenzglas mit Durchmesser $d=\qty{1.5}{\centi\metre}$ und Gesamtmasse $m=\qty{50}{\gram}$ ist unten mit Blei beschwert und schwimmt aufrecht im Wasser. Wird es ein Stück aus dem Wasser gezogen und losgelassen, schwingt es. Nehmen Sie im Bereich der Wasseroberfläche einen konstanten Querschnitt an.
\begin{center}
\CachedTikzFigure{content/resources/chapters/free_oscillations/tikz/floating_tube.tex}
\end{center}
\begin{enumerate}
\item Ermitteln Sie die rückstellende Kraft $F_\mathrm R(t)$ aus der Änderung des Auftriebs.
\item Stellen Sie die Bewegungsgleichung ohne Reibung auf und begründen Sie, warum die Schwingung harmonisch ist.
\item Berechnen Sie $\omega_0$, $f_0$ und $T_0$.
\end{enumerate}
\end{ExerciseBlock}
\UseExerciseSolution{free_oscillations_floating_test_tube}
